\documentclass{article}
\usepackage{geometry}
\geometry{margin=1.5cm, vmargin={0pt,1cm}}
\setlength{\topmargin}{-1cm}
\setlength{\headheight}{12.5pt}
\setlength{\paperheight}{29.7cm}
\setlength{\textheight}{25.3cm}

% useful packages.
\usepackage{fancyhdr}
\usepackage{layout}
\usepackage{tikz}
\usetikzlibrary{arrows.meta}

% import common macros
% % % % % % % % % % % % % % % % % % % % % % % % % % % % % % % %
% Common Macros for LaTeX
% % % % % % % % % % % % % % % % % % % % % % % % % % % % % % % %

% Useful packages
\usepackage{amsfonts}
\usepackage{amsmath}
\usepackage{amssymb}
\usepackage{amsthm}
\usepackage{enumerate}
\usepackage{graphicx}
\usepackage{multicol}
\usepackage[ruled,linesnumbered]{algorithm2e}

% Math operators and common symbols
\newcommand{\dif}{\mathrm{d}}
\newcommand{\innerProd}[1]{\left\langle #1 \right\rangle}
\newcommand{\difFrac}[2]{\frac{\dif #1}{\dif #2}}
\newcommand{\pdfFrac}[2]{\frac{\partial #1}{\partial #2}}
\newcommand{\T}{\mathrm{T}}
\newcommand{\norm}[1]{\left\| #1 \right\|}
\newcommand{\abs}[1]{\left| #1 \right|}

% Vector and Matrix shortcuts
\newcommand{\vecTwo}[2]{
  \begin{bmatrix}
    #1 \\
    #2
\end{bmatrix}}
\newcommand{\vecThree}[3]{
  \begin{bmatrix}
    #1 \\
    #2 \\
    #3
\end{bmatrix}}
\newcommand{\vecTwoH}[2]{
  \begin{bmatrix}
    #1 & #2
  \end{bmatrix}
}
\newcommand{\vecThreeH}[3]{
  \begin{bmatrix}
    #1 & #2 & #3
  \end{bmatrix}
}

% Common fields
\newcommand{\R}{\mathbb{R}}
\newcommand{\C}{\mathbb{C}}
\newcommand{\Z}{\mathbb{Z}}
\newcommand{\N}{\mathbb{N}}

% Solutions and proofs
\newenvironment{solution}{
\begin{proof}[解]}{
\end{proof}}

% Utility to get environment variables (requires lualatex)
\newcommand{\getEnv}[2]{\directlua{tex.print(os.getenv("#1") or "#2")}}


% Name and uID
\newcommand{\name}{\getEnv{NAME}{NAME}}
\newcommand{\uid}{\getEnv{UID}{UID}}
\newcommand{\course}{Complex Analysis}
\newcommand{\hwNum}{11}

\begin{document}

\pagestyle{fancy}
\fancyhead{}
\lhead{\zhstr{\name} (\uid)}
\chead{\course\ \#\hwNum}
\rhead{\today}

\section{Exercise 4.97}

Let $c \notin D$ be a singularity of an analytic function $f : D \to \C$ with
$D \subset \C$ open. Show that the following statements are equivalent:
\begin{enumerate}[(a)]
  \item $c$ is a removable singularity of $f$.
  \item There exists $r > 0$ such that $\dot{U}_r(c) \subseteq D$ and $f$ is bounded
    over $\dot{U}_r(c)$.
  \item The limit $\lim_{z \to c} f(z)$ exists.
  \item $\lim_{z \to c} (z - c)f(z) = 0$.
\end{enumerate}

\begin{solution}
  We prove the equivalence by showing (a) $\Rightarrow$ (b) $\Rightarrow$ (c) $\Rightarrow$ (d) $\Rightarrow$ (a).

  \medskip
  \noindent\textbf{(a) $\Rightarrow$ (b):}
  By Definition~4.92, $c$ is a removable singularity iff there exists an analytic
  continuation $F : D \cup \{c\} \to \C$ such that $F|_D = f$.
  Since $F$ is analytic at $c$, it is continuous at $c$, hence bounded in some
  neighborhood $U_r(c)$. Thus $f = F|_{\dot{U}_r(c)}$ is bounded on $\dot{U}_r(c)$.

  \medskip
  \noindent\textbf{(b) $\Rightarrow$ (c):}
  This is Theorem~4.93 (Riemann removability condition). If $f$ is bounded on
  $\dot{U}_r(c)$, then define
  \[
    h(z) =
    \begin{cases}
      (z - c)^2 f(z), & z \ne c,\\
      0, & z = c.
    \end{cases}
  \]
  Then $h$ is analytic on $U_r(c)$ with $h(c) = h'(c) = 0$. Expanding $h$ as a
  power series and factoring out $(z-c)^2$ gives an analytic continuation of $f$
  to $c$, which implies $\lim_{z \to c} f(z)$ exists and equals $F(c)$.

  \medskip
  \noindent\textbf{(c) $\Rightarrow$ (d):}
  If $\lim_{z \to c} f(z) = L \in \C$, then
  \[
    \lim_{z \to c} (z - c)f(z) = \Bigl(\lim_{z \to c} (z-c)\Bigr) \cdot
    \Bigl(\lim_{z \to c} f(z)\Bigr) = 0 \cdot L = 0.
  \]

  \medskip
  \noindent\textbf{(d) $\Rightarrow$ (a):}
  Assume $\lim_{z \to c} (z-c)f(z) = 0$. Define
  \[
    h(z) =
    \begin{cases}
      (z - c)^2 f(z), & z \ne c,\\
      0, & z = c.
    \end{cases}
  \]
  Then
  \[
    \lim_{z \to c} \frac{h(z) - h(c)}{z - c} = \lim_{z \to c} (z-c)f(z) = 0,
  \]
  so $h'(c) = 0$ and $h$ is analytic on $U_r(c)$. By Theorem~4.1, $h$ has a
  power series expansion $h(z) = a_2(z-c)^2 + a_3(z-c)^3 + \cdots$. Then for
  $z \ne c$, $f(z) = h(z)/(z-c)^2 = a_2 + a_3(z-c) + \cdots$, which defines
  an analytic extension of $f$ to $c$. Hence $c$ is removable by Definition~4.92.
\end{solution}

\section{Exercise 4.98}

Is $0$ a removable singularity of $g(z) = \sin \frac{1}{z}$? Why?

\begin{solution}
  No, $0$ is not a removable singularity of $g(z) = \sin \frac{1}{z}$.

  By Exercise~4.97(b), if $0$ were removable, then $g$ would be bounded in some
  punctured neighborhood $\dot{U}_r(0)$. However, for any $r > 0$, set
  $z_n = \frac{2}{(4n+1)\pi} \in \dot{U}_r(0)$ for sufficiently large $n$.
  Then $\frac{1}{z_n} = \frac{(4n+1)\pi}{2}$, so
  $\sin\frac{1}{z_n} = \sin\frac{(4n+1)\pi}{2} = 1$.
  Meanwhile, set $w_n = \frac{1}{2n\pi} \in \dot{U}_r(0)$, then
  $\sin\frac{1}{w_n} = \sin(2n\pi) = 0$.
  Thus $g$ is not bounded (and has no limit) as $z \to 0$. By the equivalence
  of (a) and (b) in Exercise~4.97, $0$ is not removable.

  In fact, $0$ is an essential singularity of $g$, as can be seen from the
  Laurent expansion $g(z) = \sum_{n=0}^{\infty} \frac{(-1)^n}{(2n+1)!} z^{-(2n+1)}$,
  which has infinitely many negative powers.
\end{solution}

\section{Exercise 4.99}

Let $c \notin D$ be a singularity of an analytic function $f : D \to \C$ with
$D \subset \C$ open. Show that $c$ is a removable singularity of $f$ if there
exist positive constants $r$, $A$, and $\epsilon$ such that
\[
  \forall z \in \dot{U}_r(c), \quad |f(z)| \le A |z - c|^{-1 + \epsilon}.
\]

\begin{solution}
  Using the given bound, for any $z \in \dot{U}_r(c)$,
  \[
    |(z - c)f(z)| = |z - c| \cdot |f(z)|
    \le |z - c| \cdot A |z - c|^{-1 + \epsilon}
    = A |z - c|^{\epsilon}.
  \]
  Since $\epsilon > 0$, we have $\lim_{z \to c} |z - c|^{\epsilon} = 0$. Hence
  \[
    \lim_{z \to c} (z - c)f(z) = 0.
  \]
  By Exercise~4.97(d), $c$ is a removable singularity of $f$.
\end{solution}

\section{Exercise 4.108}

Let $c \notin D$ be a singularity of an analytic function $f : D \to \C$ with
$D \subset \C$ open. Show that the following statements are equivalent:
\begin{enumerate}[(a)]
  \item $c$ is a pole of $f$ of order $k \in \N$.
  \item There exist $r > 0$ and an analytic function $h : U_r(c) \to \C$ with
    $h(c) \ne 0$ such that $\dot{U}_r(c) \subseteq D$ and
    $\forall z \in \dot{U}_r(c), \; f(z) = \frac{h(z)}{(z-c)^k}$.
  \item There exist $r > 0$ and an analytic function $g : U_r(c) \to \C$ not
    vanishing in $\dot{U}_r(c)$ such that $g$ has a zero of multiplicity $k$
    at $c$ and $f = \frac{1}{g}$ over $\dot{U}_r(c)$.
  \item There exist $r > 0$ and constants $M_1, M_2 > 0$ such that
    $\forall z \in \dot{U}_r(c), \; M_1 |z-c|^{-k} \le |f(z)| \le M_2 |z-c|^{-k}$.
\end{enumerate}

\begin{solution}
  We prove the equivalence cyclically.

  \medskip
  \noindent\textbf{(a) $\Rightarrow$ (b):}
  By Definition~4.103 and Definition~4.106, $c$ is a pole of order $k$ means
  $\operatorname{ord}(f; c) = -k$. Thus $h(z) := (z-c)^k f(z)$ has a removable
  singularity at $c$ with $h(c) \ne 0$. Extending $h$ analytically to $U_r(c)$
  gives the desired representation $f(z) = h(z)/(z-c)^k$.

  \medskip
  \noindent\textbf{(b) $\Rightarrow$ (c):}
  Given $f(z) = h(z)/(z-c)^k$ with $h$ analytic and $h(c) \ne 0$, define
  $g(z) = (z-c)^k/h(z)$. Then $g$ is analytic on some $U_r(c)$ (shrinking $r$
  if needed so $h(z) \ne 0$), $g(c) = 0$, and $g^{(k)}(c) \ne 0$ (since
  $h(c) \ne 0$), so $g$ has a zero of multiplicity $k$ at $c$. Moreover,
  $g(z) \ne 0$ for $z \in \dot{U}_r(c)$, and $f = 1/g$.

  \medskip
  \noindent\textbf{(c) $\Rightarrow$ (d):}
  Since $g$ has a zero of multiplicity $k$ at $c$, we can write
  $g(z) = (z-c)^k \varphi(z)$ with $\varphi$ analytic and $\varphi(c) \ne 0$.
  Then $f(z) = 1/g(z) = 1/[(z-c)^k \varphi(z)]$. Since $\varphi$ is continuous
  and nonzero at $c$, there exist $m_1, m_2 > 0$ such that
  $m_1 \le |\varphi(z)| \le m_2$ for $z$ near $c$. Setting $M_1 = 1/m_2$ and
  $M_2 = 1/m_1$ yields $M_1 |z-c|^{-k} \le |f(z)| \le M_2 |z-c|^{-k}$.

  \medskip
  \noindent\textbf{(d) $\Rightarrow$ (a):}
  From the inequality, $|(z-c)^k f(z)| \ge M_1 > 0$ and
  $|(z-c)^{k-1} f(z)| \le M_2 |z-c|^{-1}$. By the lower bound, $(z-c)^k f(z)$
  has a removable singularity at $c$ (it is bounded below, and together with the
  upper bound we show it actually has a non-zero limit). Specifically,
  $(z-c)^k f(z)$ is bounded near $c$ (by $M_2$ times a factor), hence by
  Theorem~4.93 (Exercise~4.97(b)$\Rightarrow$(a)) it has a removable singularity.
  Its limit at $c$ is non-zero because $|(z-c)^k f(z)| \ge M_1$.
  Meanwhile, $(z-c)^{k-1} f(z) \to \infty$ as $z \to c$ (since
  $|(z-c)^{k-1} f(z)| \ge M_1 |z-c|^{-1} \to \infty$), so $k$ is minimal.
  Hence $\operatorname{ord}(f; c) = -k$, i.e., $c$ is a pole of order $k$.
\end{solution}

\section{Exercise 4.114}

Find all isolated singularities of
\[
  f(z) = \frac{(z-1)^2(z+3)}{1 - \sin\frac{\pi z}{2}}
\]
and determine for each one its type.

\begin{solution}
  The singularities occur when the denominator vanishes:
  \[
    1 - \sin\frac{\pi z}{2} = 0 \iff \sin\frac{\pi z}{2} = 1
    \iff \frac{\pi z}{2} = \frac{\pi}{2} + 2\pi n,\; n \in \Z
    \iff z = 1 + 4n,\; n \in \Z.
  \]

  Let $z_n = 1 + 4n$ for $n \in \Z$. We analyze each singularity.

  \medskip
  \noindent\textbf{At $z = 1$ (i.e., $n = 0$):}
  Compute the Taylor expansion of $\sin\frac{\pi z}{2}$ at $z = 1$:
  \begin{align*}
    g(z) &:= 1 - \sin\frac{\pi z}{2}, \\
    g(1) &= 1 - \sin\frac{\pi}{2} = 0,\\
    g'(z) &= -\frac{\pi}{2} \cos\frac{\pi z}{2}, \quad
    g'(1) = -\frac{\pi}{2} \cos\frac{\pi}{2} = 0,\\
    g''(z) &= \frac{\pi^2}{4} \sin\frac{\pi z}{2}, \quad
    g''(1) = \frac{\pi^2}{4} \sin\frac{\pi}{2} = \frac{\pi^2}{4} \ne 0.
  \end{align*}
  Thus $g$ has a zero of multiplicity $2$ at $z = 1$. The numerator is
  $(z-1)^2(z+3)$, which has a zero of multiplicity $2$ at $z = 1$. Hence
  \[
    f(z) = \frac{(z-1)^2(z+3)}{g(z)} =
    \frac{(z-1)^2(z+3)}{(z-1)^2 \cdot \frac{g''(1)}{2!} + O((z-1)^3)}
    = \frac{z+3}{\frac{g''(1)}{2} + O(z-1)}.
  \]
  The factor $(z-1)^2$ cancels, so $z = 1$ is a \textbf{removable singularity}.
  Moreover, $\lim_{z \to 1} f(z) = \frac{4}{\pi^2/8} = \frac{32}{\pi^2}$.

  \medskip
  \noindent\textbf{At $z = 1 + 4n$ for $n \ne 0$:}
  Let $z_n = 1 + 4n$. Compute:
  \[
    g(z_n) = 1 - \sin\frac{\pi(1+4n)}{2} = 1 - \sin\Bigl(\frac{\pi}{2} + 2\pi n\Bigr)
    = 1 - \sin\frac{\pi}{2} = 0,
  \]
  \[
    g'(z_n) = -\frac{\pi}{2} \cos\frac{\pi(1+4n)}{2}
    = -\frac{\pi}{2} \cos\Bigl(\frac{\pi}{2} + 2\pi n\Bigr)
    = -\frac{\pi}{2} \cos\frac{\pi}{2} = 0.
  \]
  The second derivative:
  \[
    g''(z_n) = \frac{\pi^2}{4} \sin\frac{\pi(1+4n)}{2}
    = \frac{\pi^2}{4} \sin\Bigl(\frac{\pi}{2} + 2\pi n\Bigr)
    = \frac{\pi^2}{4} \ne 0.
  \]
  So for all $n \in \Z$, $g(z)$ has a zero of multiplicity $2$ at $z = 1+4n$.

  For $n \ne 0$, the numerator $(z-1)^2(z+3)$ does \emph{not} vanish at
  $z = 1+4n$ (since $1+4n \ne 1$ and $1+4n \ne -3$ for $n \ne 0, -1$).
  Check $n = -1$: $z_{-1} = -3$. Then numerator = $(-4)^2 \cdot 0 = 0$,
  with multiplicity $1$. Denominator has zero of multiplicity $2$. Hence
  $f(z) \sim \frac{C}{(z+3)^1}$ near $z = -3$, so $z = -3$ is a
  \textbf{simple pole}.

  For $n \ne 0, -1$, the numerator is non-zero and analytic at $z_n$, while
  the denominator has a zero of order $2$. Hence $f$ has a \textbf{pole of
  order $2$} at each $z_n$ with $n \in \Z \setminus \{0, -1\}$.

  \medskip
  \noindent\textbf{Summary:}
  \begin{itemize}
    \item $z = 1$: removable singularity.
    \item $z = -3$: simple pole (pole of order $1$).
    \item $z = 1 + 4n,\; n \in \Z \setminus \{0, -1\}$: poles of order $2$.
  \end{itemize}
\end{solution}

\section{Exercise 4.116}

Give an alternative proof of Theorem~4.115 by applying the Riemann removability
theorem~4.93 to the function $g_z : A_c(r, R) \setminus \{z\} \to \C$ given by
\[
  g_z(\zeta) := \frac{f(\zeta) - f(z)}{\zeta - z}.
\]
You are not supposed to use Corollary~3.76.

\begin{solution}
  Fix $z \in A_c(r, R)$. The function $g_z(\zeta) = \frac{f(\zeta)-f(z)}{\zeta-z}$
  is analytic on $A_c(r, R) \setminus \{z\}$ and continuous on $A_c(r, R)$.
  At $\zeta = z$, $g_z$ has a removable singularity because
  $\lim_{\zeta \to z} g_z(\zeta) = f'(z)$ exists. By Theorem~4.93 (Riemann
  removability), $g_z$ extends analytically to all of $A_c(r, R)$.

  Now apply the Cauchy integral theorem on the annulus. Choose $r_1, R_1$ such
  that $r < r_1 < |z-c| < R_1 < R$. Consider the cycle
  $\Gamma = \partial U_{R_1}(c) \cup (-\partial U_{r_1}(c))$ (the boundary of
  the annulus with positive orientation on the outer circle and negative on the
  inner one). By the homology form of Cauchy's theorem (Theorem~3.70), since
  $g_z$ is analytic on $A_c(r, R)$ and $\Gamma$ is homologous to $0$ in this
  domain,
  \[
    \frac{1}{2\pi i} \oint_{\Gamma} g_z(\zeta) \dd \zeta = 0.
  \]

  Expanding the integral:
  \[
    0 = \frac{1}{2\pi i} \oint_{|\zeta-c|=R_1} \frac{f(\zeta)-f(z)}{\zeta-z} \dd\zeta
    - \frac{1}{2\pi i} \oint_{|\zeta-c|=r_1} \frac{f(\zeta)-f(z)}{\zeta-z} \dd\zeta.
  \]

  Hence
  \begin{align*}
    \frac{1}{2\pi i} \oint_{|\zeta-c|=R_1} \frac{f(\zeta)}{\zeta-z} \dd\zeta
    - \frac{1}{2\pi i} \oint_{|\zeta-c|=r_1} \frac{f(\zeta)}{\zeta-z} \dd\zeta
    &= \frac{f(z)}{2\pi i} \left(
      \oint_{|\zeta-c|=R_1} \frac{\dd\zeta}{\zeta-z}
      - \oint_{|\zeta-c|=r_1} \frac{\dd\zeta}{\zeta-z}
    \right)\\
    &= \frac{f(z)}{2\pi i} \bigl(2\pi i - 0\bigr) = f(z).
  \end{align*}
  Here we used that $z$ lies inside the outer circle (winding number $1$) but
  outside the inner circle (winding number $0$). This gives exactly
  \[
    f(z) = \frac{1}{2\pi i} \oint_{|\zeta-c|=R} \frac{f(\zeta)}{\zeta-z} \dd\zeta
    - \frac{1}{2\pi i} \oint_{|\zeta-c|=r} \frac{f(\zeta)}{\zeta-z} \dd\zeta,
  \]
  which is Theorem~4.115.
\end{solution}

\section{Exercise 4.123}

Find all possible Laurent series of
\[
  f(z) := \frac{1}{z(z-1)(z-2)}.
\]

\begin{solution}
  Using partial fraction decomposition:
  \[
    f(z) = \frac{1}{z(z-1)(z-2)} = \frac{1}{2} \cdot \frac{1}{z}
    - \frac{1}{z-1} + \frac{1}{2} \cdot \frac{1}{z-2}.
  \]
  (Verify: $\frac{1}{2z} - \frac{1}{z-1} + \frac{1}{2(z-2)}
  = \frac{(z-1)(z-2) - 2z(z-2) + z(z-1)}{2z(z-1)(z-2)}
  = \frac{2}{2z(z-1)(z-2)} = \frac{1}{z(z-1)(z-2)}$.)

  The singularities are at $z = 0, 1, 2$. There are three possible annuli
  centered at $0$ (by convention, Laurent series are expanded at $c = 0$):
  \begin{enumerate}
    \item $A_0(0, 1) = \{z : 0 < |z| < 1\}$,
    \item $A_0(1, 2) = \{z : 1 < |z| < 2\}$,
    \item $A_0(2, +\infty) = \{z : |z| > 2\}$.
  \end{enumerate}

  \medskip
  \noindent\textbf{Region 1: $0 < |z| < 1$.}
  Here $|z| < 1 < 2$, so we expand each term:
  \begin{align*}
    \frac{1}{z-1} &= -\frac{1}{1-z} = -\sum_{n=0}^{\infty} z^n,\\
    \frac{1}{z-2} &= -\frac{1}{2} \cdot \frac{1}{1 - z/2}
    = -\frac{1}{2} \sum_{n=0}^{\infty} \left(\frac{z}{2}\right)^n
    = -\sum_{n=0}^{\infty} \frac{z^n}{2^{n+1}}.
  \end{align*}
  Hence
  \begin{align*}
    f(z) &= \frac{1}{2z} - \left(-\sum_{n=0}^{\infty} z^n\right)
    + \frac{1}{2}\left(-\sum_{n=0}^{\infty} \frac{z^n}{2^{n+1}}\right)\\
    &= \frac{1}{2} z^{-1} + \sum_{n=0}^{\infty} z^n
    - \sum_{n=0}^{\infty} \frac{z^n}{2^{n+2}}\\
    &= \frac{1}{2} z^{-1} + \sum_{n=0}^{\infty}
    \left(1 - \frac{1}{2^{n+2}}\right) z^n.
  \end{align*}

  \medskip
  \noindent\textbf{Region 2: $1 < |z| < 2$.}
  Here $|z| > 1$ and $|z| < 2$:
  \begin{align*}
    \frac{1}{z-1} &= \frac{1}{z} \cdot \frac{1}{1 - 1/z}
    = \frac{1}{z} \sum_{n=0}^{\infty} \left(\frac{1}{z}\right)^n
    = \sum_{n=0}^{\infty} \frac{1}{z^{n+1}} = \sum_{n=-\infty}^{-1} z^n,\\
    \frac{1}{z-2} &= -\frac{1}{2} \cdot \frac{1}{1 - z/2}
    = -\sum_{n=0}^{\infty} \frac{z^n}{2^{n+1}}.
  \end{align*}
  Hence
  \begin{align*}
    f(z) &= \frac{1}{2z} - \sum_{n=-\infty}^{-1} z^n
    + \frac{1}{2}\left(-\sum_{n=0}^{\infty} \frac{z^n}{2^{n+1}}\right)\\
    &= \frac{1}{2} z^{-1} - \sum_{n=-\infty}^{-1} z^n
    - \sum_{n=0}^{\infty} \frac{z^n}{2^{n+2}}\\
    &= -\sum_{n=-\infty}^{-2} z^n - \frac{1}{2}z^{-1}
    - \sum_{n=0}^{\infty} \frac{z^n}{2^{n+2}}.
  \end{align*}
  More compactly:
  \[
    f(z) = -\sum_{n=-\infty}^{-1} z^n - \sum_{n=0}^{\infty} \frac{z^n}{2^{n+2}}.
  \]
  (Note: the $n=-1$ term is $-z^{-1}$, combining with $\frac{1}{2}z^{-1}$ gives
  $-\frac{1}{2}z^{-1}$. Let me recheck.)

  Actually, recompute carefully:
  \begin{align*}
    f(z) &= \frac{1}{2z} - \sum_{n=1}^{\infty} \frac{1}{z^n}
    - \sum_{n=0}^{\infty} \frac{z^n}{2^{n+2}}\\
    &= \frac{1}{2}z^{-1} - z^{-1} - \sum_{n=2}^{\infty} z^{-n}
    - \sum_{n=0}^{\infty} \frac{z^n}{2^{n+2}}\\
    &= -\frac{1}{2}z^{-1} - \sum_{n=2}^{\infty} z^{-n}
    - \sum_{n=0}^{\infty} \frac{z^n}{2^{n+2}}.
  \end{align*}
  So $a_{-1} = -\frac{1}{2}$, $a_{-n} = -1$ for $n \ge 2$, and
  $a_n = -\frac{1}{2^{n+2}}$ for $n \ge 0$.

  \medskip
  \noindent\textbf{Region 3: $|z| > 2$.}
  Both $|z| > 1$ and $|z| > 2$:
  \begin{align*}
    \frac{1}{z-1} &= \sum_{n=1}^{\infty} \frac{1}{z^n},\\
    \frac{1}{z-2} &= \frac{1}{z} \cdot \frac{1}{1 - 2/z}
    = \sum_{n=1}^{\infty} \frac{2^{n-1}}{z^n}.
  \end{align*}
  Hence
  \begin{align*}
    f(z) &= \frac{1}{2z} - \sum_{n=1}^{\infty} \frac{1}{z^n}
    + \frac{1}{2} \sum_{n=1}^{\infty} \frac{2^{n-1}}{z^n}\\
    &= \frac{1}{2}z^{-1} - \sum_{n=1}^{\infty} z^{-n}
    + \sum_{n=1}^{\infty} \frac{2^{n-2}}{z^n}\\
    &= \frac{1}{2}z^{-1} + \sum_{n=1}^{\infty} (2^{n-2} - 1) z^{-n}.
  \end{align*}
  For $n = 1$: $a_{-1} = \frac{1}{2} + (2^{-1} - 1) = \frac{1}{2} - \frac{1}{2} = 0$.
  For $n \ge 2$: $a_{-n} = 2^{n-2} - 1$.

  So $f(z) = \sum_{n=2}^{\infty} (2^{n-2} - 1) z^{-n} = \sum_{n=2}^{\infty}
  \frac{2^{n-2} - 1}{z^n}$ for $|z| > 2$.
\end{solution}

\section{Exercise 4.126}

Let $f : A_c(r, R) \to \C$ be an analytic function with the Laurent
expansion (4.39). Prove that
\[
  \forall \rho \in (r, R), \quad
  \sum_{j=-\infty}^{+\infty} |a_j|^2 \rho^{2j} \le [M(\rho)]^2,
\]
where $M(\rho)$ is defined in Corollary~3.86. In particular, for any
$j \in \Z$, $|a_j| \le \rho^{-j} M(\rho)$; the equality holds for some
$j = j_0$ and $\rho = \rho_0$ if and only if $f(z) = a_{j_0}(z-c)^{j_0}$.

\begin{solution}
  From Theorem~4.124, the Laurent coefficients are given by
  \[
    a_j = \frac{1}{2\pi i} \oint_{|\zeta-c|=\rho}
    \frac{f(\zeta)}{(\zeta-c)^{j+1}} \dd\zeta, \qquad \rho \in (r, R).
  \]
  Parametrizing $\zeta = c + \rho e^{i\theta}$ for $\theta \in [0, 2\pi]$,
  \begin{align*}
    a_j &= \frac{1}{2\pi i} \int_{0}^{2\pi}
    \frac{f(c + \rho e^{i\theta})}{\rho^{j+1} e^{i(j+1)\theta}}
    \cdot i\rho e^{i\theta} \dd\theta\\
    &= \frac{\rho^{-j}}{2\pi} \int_{0}^{2\pi}
    f(c + \rho e^{i\theta}) \, e^{-ij\theta} \dd\theta.
  \end{align*}
  Thus $a_j \rho^j = \frac{1}{2\pi} \int_0^{2\pi} f(c + \rho e^{i\theta})
  e^{-ij\theta} \dd\theta$, which are precisely the Fourier coefficients of
  the function $\theta \mapsto f(c + \rho e^{i\theta})$.

  By Parseval's identity for Fourier series (or Gutzmer's inequality,
  Exercise~4.2),
  \[
    \sum_{j=-\infty}^{+\infty} |a_j|^2 \rho^{2j}
    = \frac{1}{2\pi} \int_{0}^{2\pi} |f(c + \rho e^{i\theta})|^2 \dd\theta
    \le \frac{1}{2\pi} \int_{0}^{2\pi} [M(\rho)]^2 \dd\theta
    = [M(\rho)]^2,
  \]
  where $M(\rho) = \max_{|\zeta-c|=\rho} |f(\zeta)|$ as defined in
  Corollary~3.86.

  The particular case $|a_j| \le \rho^{-j} M(\rho)$ follows immediately since
  each term in the sum is non-negative:
  \[
    |a_j|^2 \rho^{2j} \le \sum_{k=-\infty}^{+\infty} |a_k|^2 \rho^{2k}
    \le [M(\rho)]^2 \implies |a_j| \le \rho^{-j} M(\rho).
  \]

  For the equality case: if $f(z) = a_{j_0}(z-c)^{j_0}$, then by direct
  computation $M(\rho_0) = |a_{j_0}| \rho_0^{j_0}$, so
  $|a_{j_0}| = \rho_0^{-j_0} M(\rho_0)$.

  Conversely, suppose equality $|a_{j_0}| = \rho_0^{-j_0} M(\rho_0)$ holds.
  Then all other terms in the sum must be zero:
  \[
    \sum_{j \ne j_0} |a_j|^2 \rho_0^{2j}
    = [M(\rho_0)]^2 - |a_{j_0}|^2 \rho_0^{2j_0} = 0,
  \]
  so $a_j = 0$ for all $j \ne j_0$. Hence $f(z) = a_{j_0}(z-c)^{j_0}$.
  This can also be argued via the equality condition of Parseval's identity:
  $|f(c + \rho_0 e^{i\theta})| \equiv M(\rho_0)$ constant on the circle
  $|\zeta-c| = \rho_0$, which (by the maximum modulus principle) forces
  $f$ to be a monomial.
\end{solution}

\section{Exercise 4.128}

Prove Lemma~4.127.

\begin{solution}
  \textbf{Lemma~4.127.} The Laurent expansion of an analytic function
  $f : A_0(r, R) \to \C$ only contains even (or odd) powers of $z$ if $f$ is
  respectively even (or odd).

  \medskip
  \noindent\textbf{Even case:} Suppose $f(-z) = f(z)$ for all
  $z \in A_0(r, R)$. Write the Laurent expansion
  \[
    f(z) = \sum_{j=-\infty}^{\infty} a_j z^j, \qquad z \in A_0(r, R).
  \]
  Then
  \[
    f(-z) = \sum_{j=-\infty}^{\infty} a_j (-z)^j
    = \sum_{j=-\infty}^{\infty} (-1)^j a_j z^j.
  \]
  Since $f(-z) = f(z)$, the uniqueness of Laurent expansion
  (Theorem~4.124) implies $a_j = (-1)^j a_j$ for all $j \in \Z$.
  For odd $j$, $(-1)^j = -1$, so $a_j = -a_j \implies a_j = 0$.
  Hence only even powers survive.

  \medskip
  \noindent\textbf{Odd case:} Suppose $f(-z) = -f(z)$ for all
  $z \in A_0(r, R)$. Then
  \[
    f(-z) = \sum_{j=-\infty}^{\infty} (-1)^j a_j z^j
    = -f(z) = -\sum_{j=-\infty}^{\infty} a_j z^j
    = \sum_{j=-\infty}^{\infty} (-a_j) z^j.
  \]
  By uniqueness, $(-1)^j a_j = -a_j$, i.e., $a_j = (-1)^{j+1} a_j$ for all
  $j \in \Z$. For even $j$, $(-1)^{j+1} = -1$, so $a_j = -a_j \implies a_j = 0$.
  Hence only odd powers survive.
\end{solution}

\section{Exercise 4.131}

Give an alternative proof of (d) $\Rightarrow$ (a) in Exercise~4.97 by
applying Exercise~4.126.

\begin{solution}
  Recall that (d) states $\lim_{z \to c} (z-c)f(z) = 0$, and (a) states that
  $c$ is a removable singularity.

  WLOG, assume $c = 0$. Since $f$ is analytic on some punctured disk
  $\dot{U}_R(0)$, it has a Laurent expansion
  \[
    f(z) = \sum_{j=-\infty}^{\infty} a_j z^j, \qquad z \in A_0(0, R).
  \]

  Condition (d) gives $\lim_{z \to 0} z f(z) = 0$. Compute:
  \[
    z f(z) = \sum_{j=-\infty}^{\infty} a_j z^{j+1}
    = \sum_{j=-\infty}^{\infty} a_{j-1} z^j.
  \]
  For this to have limit $0$ as $z \to 0$, we need all negative-power
  coefficients to vanish.

  We prove this using Exercise~4.126. For any $\rho \in (0, R)$,
  Exercise~4.126 gives $|a_j| \le \rho^{-j} M(\rho)$.
  For $j < 0$, this becomes $|a_j| \le \rho^{-j} M(\rho)$ with $-j > 0$.

  Now, condition (d) implies that $z f(z)$ is bounded near $0$ (in fact,
  tends to $0$). Thus there exists $B > 0$ such that
  $|z f(z)| \le B$ for all $z \in \dot{U}_{\delta}(0)$ for some $\delta > 0$.
  Then for $\rho < \delta$,
  \[
    M(\rho) = \max_{|z|=\rho} |f(z)|
    \le \max_{|z|=\rho} \frac{B}{|z|} = \frac{B}{\rho}.
  \]

  Now for $j \le -2$, apply Exercise~4.126:
  \[
    |a_j| \le \rho^{-j} M(\rho) \le \rho^{-j} \cdot \frac{B}{\rho}
    = B \rho^{-j-1}.
  \]
  Since $-j-1 \ge 1$ for $j \le -2$, letting $\rho \to 0^+$ yields
  $|a_j| = 0$ for all $j \le -2$.

  For $j = -1$, we have $a_{-1} = \lim_{z \to 0} z f(z) = 0$ by condition (d).

  Therefore all negative-index coefficients vanish: $a_j = 0$ for all
  $j < 0$. By Theorem~4.130(a), $c = 0$ is a removable singularity.

  \medskip
  \noindent\textit{Remark.} This proof shows that condition (d) forces the
  principal part of the Laurent series to vanish, which is exactly the
  characterization of removable singularities via Laurent series
  (Theorem~4.130).
\end{solution}

\section{Exercise 4.134 (Laurent series as an expansion in an orthonormal basis)}

Denote by $H$ the space of all analytic functions on $A_0(r, R)$ where
$r < 1 < R$ and define a sesquilinear map $\langle\cdot, \cdot\rangle :
H \times H \to \C$,
\[
  \langle f, g \rangle := \frac{1}{2\pi i}
  \oint_{|\zeta| = 1} f(\zeta) \overline{g(\zeta)} \frac{1}{\zeta} \dd\zeta.
\]
Show that
\begin{enumerate}[(a)]
  \item $\langle\cdot, \cdot\rangle$ is an inner product on $H$ in the sense
    of Definition~C.170;
  \item the set $\{z^j : j \in \Z\}$ is an orthonormal basis of $H$;
  \item the Laurent series of any $f \in H$ is exactly the expansion of $f$ in
    terms of the orthonormal basis $\{z^j : j \in \Z\}$. In other words, if
    $f(z) = \sum_{j=-\infty}^{+\infty} a_j z^j$ is the Laurent expansion of
    $f$, then
    \[
      \forall j \in \Z, \quad a_j = \langle f, z^j \rangle
      = \frac{1}{2\pi i} \oint_{|\zeta|=1} \frac{f(\zeta)}{\zeta^{j+1}} \dd\zeta.
    \]
\end{enumerate}

\begin{solution}
  \textbf{(a) Inner product verification.}

  Recall (Definition~C.170) that $\langle\cdot,\cdot\rangle$ is an inner
  product if it is: (i) positive definite,
  (ii) conjugate symmetric, (iii) linear in the first argument.

  \begin{itemize}
    \item \textbf{Conjugate symmetry:}
    \begin{align*}
      \overline{\langle g, f \rangle}
      &= \overline{\frac{1}{2\pi i} \oint_{|\zeta|=1}
        g(\zeta) \overline{f(\zeta)} \frac{1}{\zeta} \dd\zeta}
      = -\frac{1}{2\pi i} \oint_{|\zeta|=1}
        \overline{g(\zeta)} f(\zeta) \frac{1}{\bar{\zeta}} \dd\bar{\zeta}.
    \end{align*}
    On the unit circle, $\zeta = e^{i\theta}$, $\bar{\zeta} = e^{-i\theta} = 1/\zeta$,
    $\dd\bar{\zeta} = -ie^{-i\theta} \dd\theta$, and $\dd\zeta = ie^{i\theta} \dd\theta$.
    Through the change of variables or using $\bar{\zeta} = 1/\zeta$ and
    $\dd\bar{\zeta} = -\dd\zeta/\zeta^2$, a direct computation shows
    $\overline{\langle g, f \rangle} = \langle f, g \rangle$.

    More concretely, parametrize $\zeta = e^{i\theta}$:
    \[
      \langle f, g \rangle = \frac{1}{2\pi i} \int_0^{2\pi}
      f(e^{i\theta}) \overline{g(e^{i\theta})} \frac{1}{e^{i\theta}}
      \cdot ie^{i\theta} \dd\theta
      = \frac{1}{2\pi} \int_0^{2\pi}
      f(e^{i\theta}) \overline{g(e^{i\theta})} \dd\theta.
    \]
    Then $\overline{\langle g, f \rangle} = \overline{\frac{1}{2\pi}
    \int_0^{2\pi} g(e^{i\theta}) \overline{f(e^{i\theta})} \dd\theta}
    = \frac{1}{2\pi} \int_0^{2\pi} f(e^{i\theta})
    \overline{g(e^{i\theta})} \dd\theta = \langle f, g \rangle$.

    \item \textbf{Linearity in the first argument:}
    For $\alpha, \beta \in \C$ and $f_1, f_2, g \in H$,
    \begin{align*}
      \langle \alpha f_1 + \beta f_2, g \rangle
      &= \frac{1}{2\pi} \int_0^{2\pi}
      (\alpha f_1(e^{i\theta}) + \beta f_2(e^{i\theta}))
      \overline{g(e^{i\theta})} \dd\theta\\
      &= \alpha \langle f_1, g \rangle + \beta \langle f_2, g \rangle.
    \end{align*}

    \item \textbf{Positive definiteness:}
    $\langle f, f \rangle = \frac{1}{2\pi} \int_0^{2\pi}
    |f(e^{i\theta})|^2 \dd\theta \ge 0$. If $\langle f, f \rangle = 0$,
    then $f(e^{i\theta}) = 0$ for all $\theta \in [0, 2\pi]$ (by continuity
    of $|f|^2$). Thus $f$ vanishes on the unit circle. Since the unit circle
    has an accumulation point in the domain $A_0(r, R)$, the identity theorem
    (Theorem~4.35) implies $f \equiv 0$ on $A_0(r, R)$.
  \end{itemize}
  Hence $\langle\cdot,\cdot\rangle$ is an inner product on $H$.

  \medskip
  \textbf{(b) Orthonormality of $\{z^j : j \in \Z\}$.}

  For any $j, k \in \Z$,
  \[
    \langle z^j, z^k \rangle
    = \frac{1}{2\pi i} \oint_{|\zeta|=1} \zeta^j \overline{\zeta^k}
    \frac{1}{\zeta} \dd\zeta
    = \frac{1}{2\pi i} \oint_{|\zeta|=1} \zeta^j \zeta^{-k}
    \frac{1}{\zeta} \dd\zeta
    = \frac{1}{2\pi i} \oint_{|\zeta|=1} \zeta^{j-k-1} \dd\zeta.
  \]
  (Here we used $\overline{\zeta} = 1/\zeta$ on the unit circle.)

  By the fundamental integral formula (or Example~3.16),
  \[
    \frac{1}{2\pi i} \oint_{|\zeta|=1} \zeta^{m} \dd\zeta
    = \begin{cases}
      1, & m = -1,\\
      0, & m \ne -1.
    \end{cases}
  \]
  Here $m = j-k-1$. Thus $m = -1 \iff j-k-1 = -1 \iff j = k$.
  Hence $\langle z^j, z^k \rangle = \delta_{jk}$, proving orthonormality.

  To see that $\{z^j\}$ is a \textbf{basis} (i.e., complete): any
  $f \in H$ has a Laurent expansion $f(z) = \sum_{j=-\infty}^{\infty} a_j z^j$
  converging normally on $A_0(r, R)$. The convergence is uniform on the unit
  circle, so $f$ can be approximated in the $L^2$-norm induced by the inner
  product by finite linear combinations of $\{z^j\}$. Hence $\{z^j\}$ spans
  a dense subspace; since it is orthonormal, it is an orthonormal basis.

  \medskip
  \textbf{(c) Laurent coefficients as inner products.}

  For any $j \in \Z$,
  \begin{align*}
    \langle f, z^j \rangle
    &= \frac{1}{2\pi i} \oint_{|\zeta|=1} f(\zeta) \overline{\zeta^j}
    \frac{1}{\zeta} \dd\zeta
    = \frac{1}{2\pi i} \oint_{|\zeta|=1} f(\zeta) \zeta^{-j}
    \frac{1}{\zeta} \dd\zeta\\
    &= \frac{1}{2\pi i} \oint_{|\zeta|=1}
    \frac{f(\zeta)}{\zeta^{j+1}} \dd\zeta.
  \end{align*}

  By Theorem~4.124, the Laurent coefficient $a_j$ is exactly
  \[
    a_j = \frac{1}{2\pi i} \oint_{|\zeta|=\rho}
    \frac{f(\zeta)}{\zeta^{j+1}} \dd\zeta
  \]
  for any $\rho \in (r, R)$. Taking $\rho = 1$ (which is valid since
  $r < 1 < R$) gives $a_j = \langle f, z^j \rangle$.

  Thus the Laurent series $f(z) = \sum_{j=-\infty}^{\infty} a_j z^j$ is
  precisely the expansion of $f$ in the orthonormal basis $\{z^j : j \in \Z\}$,
  with coefficients given by the inner product $a_j = \langle f, z^j \rangle$.
\end{solution}

\end{document}
